% =====================================================================
\chapter{Procesos de decisión de Markov y el LLM como política}
\label{tema:1}
% =====================================================================

Este tema introduce el lenguaje del aprendizaje por refuerzo: agente, entorno, recompensa, retorno, política y funciones de valor. La presentación sigue el capítulo~3 de Sutton y Barto~\cite{sb} y la clase~2 del curso de David Silver~\cite{silver}; la última sección traslada el formalismo a la generación de texto, como en la clase~17 de CS336~\cite{cs336}.

\section{La interfaz agente-entorno}

En cada instante $t=0,1,2,\dots$ un \emph{agente} observa el estado $S_t$ del \emph{entorno}, elige una acción $A_t$ y recibe, en el instante siguiente, una recompensa numérica $R_{t+1}$ y un nuevo estado $S_{t+1}$. La sucesión resultante
\[
S_0,\,A_0,\,R_1,\,S_1,\,A_1,\,R_2,\,S_2,\,\dots
\]
se llama \emph{trayectoria}. Cuando la interacción termina en un estado terminal tras un número finito de pasos, la trayectoria se llama \emph{episodio}.

\begin{definicion}[Proceso de decisión de Markov finito]\label{def:mdp}
Un \emph{proceso de decisión de Markov} (MDP) finito es una tupla $(\Sset,\Aset,p,\rho_0,\gamma)$ formada por un conjunto finito de estados $\Sset$, un conjunto finito de acciones $\Aset$, una \emph{dinámica}
\[
p(s',r\mid s,a)=P\bigl(S_{t+1}=s',\,R_{t+1}=r\mid S_t=s,\,A_t=a\bigr),
\]
que para cada par $(s,a)$ es una distribución sobre $\Sset\times\R$ con soporte finito, una distribución inicial $\rho_0$ sobre $\Sset$ y un \emph{factor de descuento} $\gamma\in[0,1]$. La palabra «Markov» expresa que la distribución de $(S_{t+1},R_{t+1})$ depende de la historia solo a través de $(S_t,A_t)$.
\end{definicion}

\begin{definicion}[Política]\label{def:politica}
Una \emph{política} (estacionaria y estocástica) es una aplicación que asigna a cada estado $s$ una distribución $\pi(\cdot\mid s)$ sobre $\Aset$. Una política \emph{parametrizada} $\pi_\theta$ depende de un vector de parámetros $\theta\in\R^{d}$.
\end{definicion}

Fijados un MDP y una política, la probabilidad de un prefijo de trayectoria $\tau=(s_0,a_0,r_1,s_1,\dots,s_T)$ se obtiene por la regla de la cadena (módulo~1):
\begin{equation}\label{eq:ptau}
p_\pi(\tau)=\rho_0(s_0)\prod_{t=0}^{T-1}\pi(a_t\mid s_t)\,p(s_{t+1},r_{t+1}\mid s_t,a_t).
\end{equation}
Esta factorización separa lo que controla el agente (los factores $\pi$) de lo que impone el entorno (los factores $p$), y es la base del teorema del gradiente de política del \cref{tema:2}.

\section{Recompensa y retorno}

El objetivo del agente no es maximizar la recompensa inmediata sino una cantidad acumulada~\cite[cap.~3]{sb}.

\begin{definicion}[Retorno]\label{def:retorno}
El \emph{retorno} desde el instante $t$ es
\[
G_t=R_{t+1}+\gamma R_{t+2}+\gamma^{2}R_{t+3}+\dots=\sum_{k=0}^{\infty}\gamma^{k}R_{t+k+1}.
\]
En una tarea episódica de horizonte $T$ se toma $R_k=0$ para $k>T$, de modo que la suma es finita y se admite $\gamma=1$.
\end{definicion}

\begin{proposicion}\label{prop:retorno}
\begin{enumerate}
  \item Para todo $t$, $G_t=R_{t+1}+\gamma\,G_{t+1}$.
  \item Si $\abs{R_k}\le R_{\max}$ para todo $k$ y $\gamma<1$, entonces $\abs{G_t}\le R_{\max}/(1-\gamma)$.
\end{enumerate}
\end{proposicion}

\begin{proof}
(i) Separando el primer término de la serie, $G_t=R_{t+1}+\sum_{k\ge 1}\gamma^{k}R_{t+k+1}=R_{t+1}+\gamma\sum_{j\ge 0}\gamma^{j}R_{t+j+2}=R_{t+1}+\gamma G_{t+1}$.
(ii) Por la desigualdad triangular, $\abs{G_t}\le\sum_{k\ge0}\gamma^{k}R_{\max}=R_{\max}/(1-\gamma)$, suma de una serie geométrica de razón $\gamma<1$.
\end{proof}

La recursión (i) permite calcular todos los retornos de un episodio en una sola pasada hacia atrás, empezando por $G_T=0$ (\cref{eje:1-retornos}).

\section{Funciones de valor y ecuación de Bellman}

\begin{definicion}[Funciones de valor y ventaja]\label{def:valor}
Sea $\pi$ una política. La \emph{función de valor de estado} y la \emph{función de valor de acción} son
\[
V^{\pi}(s)=\E_{\pi}\bigl[G_t\mid S_t=s\bigr],\qquad Q^{\pi}(s,a)=\E_{\pi}\bigl[G_t\mid S_t=s,\,A_t=a\bigr],
\]
donde $\E_\pi$ indica que las acciones futuras se eligen con $\pi$. La \emph{función ventaja} es $A^{\pi}(s,a)=Q^{\pi}(s,a)-V^{\pi}(s)$.
\end{definicion}

La ventaja mide cuánto mejor (o peor) que la media de la propia política es elegir $a$ en $s$ y continuar después con $\pi$. Es la cantidad que pondera los gradientes de todos los métodos de los temas~2, 3 y~6.

\begin{teorema}[Ecuación de Bellman para $V^{\pi}$]\label{teo:bellman}
Para todo estado $s$,
\[
V^{\pi}(s)=\sum_{a}\pi(a\mid s)\,Q^{\pi}(s,a),\qquad Q^{\pi}(s,a)=\sum_{s',r}p(s',r\mid s,a)\bigl[r+\gamma V^{\pi}(s')\bigr].
\]
En particular, $V^{\pi}(s)=\sum_{a}\pi(a\mid s)\sum_{s',r}p(s',r\mid s,a)\bigl[r+\gamma V^{\pi}(s')\bigr]$.
\end{teorema}

\begin{proof}
La primera igualdad es la ley de la esperanza total condicionando a $A_t$: $\E_\pi[G_t\mid S_t=s]=\sum_a P(A_t=a\mid S_t=s)\,\E_\pi[G_t\mid S_t=s,A_t=a]$, y $P(A_t=a\mid S_t=s)=\pi(a\mid s)$. Para la segunda, por la \cref{prop:retorno}, $G_t=R_{t+1}+\gamma G_{t+1}$. Condicionando además a $(S_{t+1},R_{t+1})=(s',r)$, que ocurre con probabilidad $p(s',r\mid s,a)$,
\[
\E_\pi[G_t\mid s,a,s',r]=r+\gamma\,\E_\pi[G_{t+1}\mid S_t=s,A_t=a,S_{t+1}=s',R_{t+1}=r]=r+\gamma V^{\pi}(s'),
\]
donde la última igualdad usa la propiedad de Markov: dado $S_{t+1}=s'$, la evolución futura no depende de la historia anterior. Sumando sobre $(s',r)$ se obtiene la expresión de $Q^{\pi}$.
\end{proof}

\begin{corolario}\label{cor:ventaja-media}
Para todo estado $s$, $\displaystyle\sum_{a}\pi(a\mid s)\,A^{\pi}(s,a)=0$.
\end{corolario}

\begin{proof}
Por el \cref{teo:bellman}, $\sum_a\pi(a\mid s)A^{\pi}(s,a)=\sum_a\pi(a\mid s)Q^{\pi}(s,a)-V^{\pi}(s)\sum_a\pi(a\mid s)=V^{\pi}(s)-V^{\pi}(s)=0$.
\end{proof}

La ecuación de Bellman es un sistema lineal de $\abs{\Sset}$ ecuaciones con $\abs{\Sset}$ incógnitas. Para $\gamma<1$ tiene solución única (\cref{eje:1-contraccion}); el capítulo~4 de~\cite{sb} la resuelve de forma iterativa.

\section{El modelo de lenguaje como política}

La generación autorregresiva encaja en la \cref{def:mdp} con una dinámica especialmente sencilla~\cite{cs336}.

\begin{definicion}[MDP de tokens]\label{def:mdp-tokens}
Sean $\V$ un vocabulario con un símbolo de fin y $x$ una instrucción. El \emph{MDP de tokens} asociado a $x$ tiene como estados las secuencias $s_t=(x,y_1,\dots,y_t)$, como acciones los tokens $a_t=y_{t+1}\in\V$, una transición determinista $s_{t+1}=(s_t,a_t)$ (concatenación), estado inicial $s_0=x$ y estados terminales los que acaban en el símbolo de fin o alcanzan una longitud máxima. La recompensa es nula salvo en la transición final, en la que vale $r(x,y)$ para la respuesta completa $y=y_{1:T}$, y se toma $\gamma=1$.
\end{definicion}

\begin{proposicion}\label{prop:mdp-tokens}
En el MDP de tokens, con una política $\pi_\theta(y_{t+1}\mid x,y_{1:t})$:
\begin{enumerate}
  \item la probabilidad de una trayectoria coincide con la de la respuesta, $p_{\pi_\theta}(\tau)=\pi_\theta(y\mid x)=\prod_{t=1}^{T}\pi_\theta(y_t\mid x,y_{<t})$;
  \item el retorno desde cualquier instante es $G_t=r(x,y)$, y $V^{\pi_\theta}(x)=\E_{y\sim\pi_\theta(\cdot\mid x)}[r(x,y)]$;
  \item para un prefijo $s_t$ y un token $a$, $Q^{\pi_\theta}(s_t,a)=V^{\pi_\theta}(s_t a)$, donde $s_t a$ denota la concatenación.
\end{enumerate}
\end{proposicion}

\begin{proof}
(i) En \eqref{eq:ptau}, $\rho_0$ es la masa puntual en $x$ y cada factor de transición vale~1, porque la transición es determinista; quedan los factores de la política, cuyo producto es $\pi_\theta(y\mid x)$ por la regla de la cadena para secuencias del módulo~1, aplicada de forma condicional a $x$. (ii) Con $\gamma=1$ y una única recompensa no nula en el último paso, la suma de la \cref{def:retorno} se reduce a $r(x,y)$; tomando esperanzas desde $s_0=x$ se obtiene $V^{\pi_\theta}(x)$. (iii) Por el \cref{teo:bellman}, con recompensa intermedia nula y transición determinista a $s_ta$, $Q^{\pi_\theta}(s_t,a)=0+1\cdot V^{\pi_\theta}(s_ta)$.
\end{proof}

\begin{table}[h]
\centering\small
\begin{tabular}{@{}ll@{}}
\toprule
Elemento del MDP & En un modelo de lenguaje \\
\midrule
Estado $s_t$ & instrucción $x$ más los tokens generados $y_{1:t}$ \\
Acción $a_t$ & siguiente token $y_{t+1}$; $\abs{\Aset}=\abs{\V}$, del orden de $10^{4}$--$10^{5}$ \\
Transición & determinista y conocida: concatenar el token \\
Política $\pi_\theta(a\mid s)$ & softmax de los logits del Transformer (módulo~14) \\
Recompensa & normalmente solo al final: $r(x,y)$ \\
Episodio & una respuesta completa, hasta el símbolo de fin \\
\bottomrule
\end{tabular}
\caption{Correspondencia entre los elementos de un MDP y la generación de texto.}
\label{tab:mdp-llm}
\end{table}

\begin{figure}[h]
\centering
\begin{tikzpicture}[font=\small,>={Stealth[length=2mm]},
  caja/.style={draw,thick,rounded corners=2pt,minimum height=9mm,align=center,inner sep=4pt}]
  \node[caja] (s) {estado $s_t$\\$x,\,y_{1:t}$};
  \node[caja,right=1.3cm of s] (p) {política $\pi_\theta$\\Transformer $+$ softmax};
  \node[caja,right=1.3cm of p] (a) {acción $a_t$\\token $y_{t+1}$};
  \node[caja,below=1.4cm of p] (r) {recompensa $r(x,y)$\\(solo al terminar)};
  \draw[->,thick] (s) -- (p);
  \draw[->,thick] (p) -- (a);
  \draw[->,thick] (a.north) -- ++(0,0.7) -| node[pos=0.25,above] {concatenar: $s_{t+1}=s_t\,a_t$} (s.north);
  \draw[->,thick,dashed] (a.south) |- node[pos=0.3,right] {¿fin?} (r.east);
  \draw[->,thick,dashed] (r.north) -- node[right] {señal de aprendizaje} (p.south);
\end{tikzpicture}
\caption{El bucle agente-entorno en la generación de texto. La transición es determinista; la única señal externa es la recompensa final.}
\label{fig:bucle}
\end{figure}

\begin{observacion}[Bandido contextual]\label{obs:bandido}
Como la transición es conocida y la recompensa llega solo al final, la respuesta entera $y$ puede verse como una única acción tomada en el contexto $x$ con probabilidad $\pi_\theta(y\mid x)$. Este caso particular de MDP de un solo paso se llama \emph{bandido contextual}. Los métodos de preferencias de los temas~4 y~5 trabajan a este nivel; PPO (\cref{tema:3}) trabaja token a token con ventajas $A^{\pi}(s_t,a_t)$. El número de acciones del bandido es enorme: con $\abs{\V}=5\cdot 10^{4}$ y $T=100$ hay $\abs{\V}^{T}\approx 10^{470}$ respuestas, por lo que ninguna suma sobre $y$ puede calcularse explícitamente.
\end{observacion}

\begin{ejemplo}[Valores y ventajas en un árbol de respuestas]\label{ej:arbol}
Sea $\V=\{\texttt{a},\texttt{b}\}$ y respuestas de exactamente dos tokens. La política elige el primer token \texttt{a} con probabilidad $0{,}6$; tras \texttt{a}, elige \texttt{a} con probabilidad $0{,}3$, y tras \texttt{b}, con probabilidad $0{,}5$. La recompensa vale~1 si la respuesta contiene exactamente una \texttt{a} y~0 en otro caso. Por la \cref{prop:mdp-tokens},
\[
V(x\texttt{a})=0{,}3\cdot 0+0{,}7\cdot 1=0{,}7,\qquad V(x\texttt{b})=0{,}5\cdot 1+0{,}5\cdot 0=0{,}5,
\]
\[
V(x)=0{,}6\cdot 0{,}7+0{,}4\cdot 0{,}5=0{,}62.
\]
Las ventajas del primer token son $A(x,\texttt{a})=0{,}7-0{,}62=0{,}08$ y $A(x,\texttt{b})=0{,}5-0{,}62=-0{,}12$, y su media ponderada es $0{,}6\cdot 0{,}08+0{,}4\cdot(-0{,}12)=0$, como exige el \cref{cor:ventaja-media}. Aunque ninguna recompensa se asigna al primer token, la ventaja le atribuye una parte del mérito final: este es el problema de \emph{asignación de crédito} que resuelven los críticos del \cref{tema:2}.
\end{ejemplo}

\begin{lecturas}
\item Sutton y Barto~\cite{sb}, cap.~3, «Finite Markov Decision Processes»: interfaz agente-entorno, retornos, políticas y funciones de valor.
\item D.~Silver~\cite{silver}, clase~2, «Markov Decision Processes».
\item CS336~\cite{cs336}, clase~17: formulación del modelo de lenguaje como política.
\item Levine, CS285~\cite{cs285}, clase~4, «Introduction to Reinforcement Learning».
\end{lecturas}

\section*{Ejercicios del tema 1}
\addcontentsline{toc}{section}{Ejercicios del tema 1}

\begin{ejercicio}{1}{eje:1-chatbot}
Un asistente conversacional recibe una nota numérica al final de cada respuesta. Identifique el estado, la acción, la transición, la recompensa y el episodio del MDP correspondiente, y explique por qué puede tratarse como un bandido contextual.
\end{ejercicio}

\begin{ejercicio}{1}{eje:1-contar}
Con $\abs{\V}=5\cdot 10^{4}$ y respuestas de longitud fija $T=100$, calcule $\log_{10}$ del número de respuestas posibles. ¿Qué consecuencia tiene para el cálculo exacto de $V^{\pi}(x)=\sum_y\pi(y\mid x)\,r(x,y)$?
\end{ejercicio}

\begin{ejercicio}{2}{eje:1-retornos}
En un episodio con $\gamma=0{,}5$ y $T=5$ se reciben las recompensas $R_1=-1$, $R_2=2$, $R_3=6$, $R_4=3$ y $R_5=2$. Calcule $G_0,G_1,\dots,G_5$ con la recursión de la \cref{prop:retorno}.
\end{ejercicio}

\begin{ejercicio}{2}{eje:1-arbol}
Repita el \cref{ej:arbol} con la política siguiente: primer token \texttt{a} con probabilidad $0{,}5$; tras \texttt{a}, \texttt{a} con probabilidad $0{,}2$; tras \texttt{b}, \texttt{a} con probabilidad $0{,}6$. Las recompensas son $r(\texttt{aa})=0$, $r(\texttt{ab})=1$, $r(\texttt{ba})=2$ y $r(\texttt{bb})=0$. Calcule $V(x)$, $V(x\texttt{a})$, $V(x\texttt{b})$, las ventajas del primer token y las del segundo token tras \texttt{b}. Compruebe el \cref{cor:ventaja-media} en ambos estados.
\end{ejercicio}

\begin{ejercicio}{2}{eje:1-cs285}
Sea un MDP con un estado no terminal $s_1$ y un estado terminal $s_F$. En $s_1$ la acción $a_1$ da recompensa~1 y mantiene el estado; la acción $a_2$ da recompensa~0 y lleva a $s_F$. La política es $\pi_\theta(a_1\mid s_1)=\theta$, $\pi_\theta(a_2\mid s_1)=1-\theta$, con $0<\theta<1$. (a) Plantee la ecuación de Bellman para $V(s_1)$ con descuento $\gamma$ y resuélvala. (b) Evalúe para $\theta=0{,}75$ con $\gamma=1$ y con $\gamma=0{,}9$. (c) Compruebe el caso $\gamma=1$ calculando directamente $\E[R(\tau)]$ como esperanza del número de veces que se elige $a_1$.
\fuente{adaptado de Levine, CS285, otoño de 2023, tarea~2 «Policy Gradients», §8, problema~1~\cite{cs285hw2}.}
\end{ejercicio}

\begin{ejercicio}{3}{eje:1-contraccion}
Sea $\gamma<1$ y sea $T^{\pi}$ el operador que asigna a cada $v\in\R^{\abs{\Sset}}$ el vector
$(T^{\pi}v)(s)=\sum_a\pi(a\mid s)\sum_{s',r}p(s',r\mid s,a)\,[r+\gamma v(s')]$. (a) Demuestre que $\norm{T^{\pi}u-T^{\pi}v}_\infty\le\gamma\norm{u-v}_\infty$. (b) Deduzca que la ecuación de Bellman tiene a lo sumo una solución. (c) Demuestre que la tiene, observando que es un sistema lineal $(I-\gamma P^{\pi})v=r^{\pi}$ y que $I-\gamma P^{\pi}$ es invertible.
\fuente{adaptado de Sutton y Barto~\cite{sb}, cap.~3, y de Silver~\cite{silver}, clase~2.}
\end{ejercicio}
